\section{CARL: Chat Assisted Rubric Learning}

CARL (Chat Assisted Rubric Learning) is a conversational framework for
translating domain expertise into fully specified reinforcement learning
systems. Rather than requiring the user to possess deep RL knowledge
upfront, CARL structures the design process as a guided dialogue across
nine sequential steps: problem definition, reference ingestion, literature
search, parameter exploration, state space design, action space definition,
reward shaping, initial testing, and full training. At each step, CARL
elicits the decisions required for that stage, suggests sensible defaults
drawn from domain context and literature, and summarises the outcomes
before advancing. The result is a complete, executable RL codebase whose
design choices are fully traceable to the domain expert's stated objectives
and constraints.

The nine steps are:

\begin{enumerate}
    \item[\textbf{1.}] \textbf{Problem Definition} --- Elicit domain objectives, dynamics, and episode structure.
    \item[\textbf{2.}] \textbf{References} --- Ingest user-provided documents and extract key domain knowledge.
    \item[\textbf{2.5.}] \textbf{Literature Search} --- Query Semantic Scholar for relevant prior work.
    \item[\textbf{3.}] \textbf{Parameter Exploration} --- Confirm domain parameter ranges for randomisation.
    \item[\textbf{4.}] \textbf{Information Proposal} --- Design state variables, features, and normalisation.
    \item[\textbf{5.}] \textbf{Action Definition} --- Define action space, bounds, and algorithm family.
    \item[\textbf{6.}] \textbf{Reward Definition} --- Design reward function, shaping, and constraints.
    \item[\textbf{7.}] \textbf{ML/RL Parameters} --- Select and confirm hyperparameters and network architecture.
    \item[\textbf{8.}] \textbf{Initial Testing} --- Generate environment, run dry-runs and baseline checks.
    \item[\textbf{9.}] \textbf{Launch} --- Execute full training pipeline and produce evaluation plots.
\end{enumerate}

The system is implemented as a multi-turn LLM agent powered by Claude via
AWS Bedrock. Steps 1--7 are conversational: the user and CARL alternate
turns to converge on a formulation, aided by Trillian, a secondary
persona that proposes user responses to accelerate the dialogue. Steps 8
and 9 are autonomous: CARL generates, executes, and iteratively debugs
all code---environment, agent, training loop, and evaluation---without
further input. This separation of concerns ensures that domain knowledge
is captured precisely during the interactive phase, while the technical
implementation is handled reliably during the autonomous phase.
